\begingroup
\renewcommand{\arraystretch}{1.45}
\Needspace{28\baselineskip}
\InterpolationArticle{33}
Zweitens sei $X$ eine Funktion der Form
\[
\begin{aligned}
&\beta'\sin x+\beta''\sin2x+\beta'''\sin3x+\text{usw.}+\beta^n\sin nx
\end{aligned}
\]
somit muss $\mu=n$ sein. Hier wird nach Artikel 30 der Koeffizient von $A$ in der zweiten Formel des Artikels 16 für $T$
\[
\begin{aligned}
&=\frac{1}{\mu\sin\mu a}\times\frac{\sin t}{\sin a}(\sin\mu a+2\sin(\mu-1)a\cos t\\[5pt]
&\qquad{}+2\sin(\mu-2)a\cos2t+\text{usw.}\\[5pt]
&\qquad+2\sin a\cos(\mu-1)t)\\[5pt]
&=\frac{1}{\mu\sin\mu a}(2\cos(\mu-1)a\sin t+2\cos(\mu-2)a\sin2t\\[5pt]
&\qquad{}+2\cos(\mu-3)a\sin3t+\text{usw.}\\[5pt]
&\qquad+2\cos a\sin(\mu-1)t+\sin\mu t)
\end{aligned}
\]
\Needspace{37\baselineskip}
Ganz ähnliche Ausdrücke ergeben sich für die Koeffizienten von $B,C,D$ usw., indem man lediglich den Bogen $a$ in $b,c,d$ usw. ändert. Setzt man daher
\[
\begin{aligned}
&\zeta'=\frac{2}{\mu\sin\mu a}\{A\cos(\mu-1)a+B\cos(\mu-1)b\\[5pt]
&\qquad{}+C\cos(\mu-1)c+D\cos(\mu-1)d+\text{usw.}\}\\[5pt]
&\zeta''=\frac{2}{\mu\sin\mu a}\{A\cos(\mu-2)a+B\cos(\mu-2)b\\[5pt]
&\qquad{}+C\cos(\mu-2)c+D\cos(\mu-2)d+\text{usw.}\}\\[5pt]
&\zeta'''=\frac{2}{\mu\sin\mu a}\{A\cos(\mu-3)a+B\cos(\mu-3)b\\[5pt]
&\qquad{}+C\cos(\mu-3)c+D\cos(\mu-3)d+\text{usw.}\}\\[5pt]
&\text{usw.}\\[5pt]
&\zeta^{\mu-1}=\frac{2}{\mu\sin\mu a}\{A\cos a+B\cos b+C\cos c+D\cos d+\text{usw.}\}\\[5pt]
&\zeta^\mu=\frac{1}{\mu\sin\mu a}\{A+B+C+D+\text{usw.}\}
\end{aligned}
\]
so wird
\[
\begin{aligned}
&T=\zeta'\sin t+\zeta''\sin2t+\zeta'''\sin3t+\text{usw.}+\zeta^\mu\sin\mu t
\end{aligned}
\]
Da diese Formel allgemein für jeden beliebigen Wert von $t$ gilt, wird sie notwendig
\SourceLeaf{325}
mit $X$ identisch sein. Ändert man $t$ in $x$, so werden daher die Koeffizienten $\zeta',\zeta'',\zeta'''$ usw. jeweils $\beta',\beta'',\beta'''$ usw. gleich sein.

\Needspace{41\baselineskip}
\InterpolationArticle{34}
Wenn wie in Artikel 31 alle gegebenen Werte durch eine solche Funktion
\[
\begin{aligned}
&X'={}\gamma+\gamma'\cos x+\gamma''\cos 2x+\text{usw.}+\gamma^{m}\cos mx\\[5pt]
&+\delta'\sin x+\delta''\sin 2x+\text{usw.}+\delta^{m}\sin mx
\end{aligned}
\]
erfüllt worden sind, wobei $m=\tfrac{1}{2}\mu-\tfrac{1}{2}$ oder $=\tfrac{1}{2}\mu$ ist, je nachdem, ob $\mu$ ungerade oder gerade ist, liefert der Vergleich dieser Funktion mit der folgenden
\[
\begin{aligned}
&X''=\zeta'\sin x+\zeta''\sin 2x+\zeta'''\sin 3x+\text{usw.}+\zeta^{\mu}\sin\mu x
\end{aligned}
\]
die wir im vorigen Artikel gewonnen haben, diese Gleichungen:
\[
\begin{aligned}
&\zeta^{\mu}=\gamma\operatorname{cosec}\mu a\\[5pt]
&\zeta'=\delta'+\gamma'\operatorname{cotg}\mu a,\\[5pt]
&\zeta^{\mu-1}=\gamma'\operatorname{cosec}\mu a\\[5pt]
&\zeta''=\delta''+\gamma''\operatorname{cotg}\mu a,\\[5pt]
&\zeta^{\mu-2}=\gamma''\operatorname{cosec}\mu a\\[5pt]
&\zeta'''=\delta'''+\gamma'''\operatorname{cotg}\mu a,\\[5pt]
&\zeta^{\mu-3}=\gamma'''\operatorname{cosec}\mu a\\[5pt]
&\text{usw., bis zu}\\[5pt]
&(\tfrac{1}{2})\zeta^{m}=\delta^{m}+\gamma^{m}\operatorname{cotg}\mu a,\\[5pt]
&(\tfrac{1}{2})\zeta^{\mu-m}=\gamma^{m}\operatorname{cosec}\mu a
\end{aligned}
\]
\Needspace{40\baselineskip}
Hier gilt der Faktor $\tfrac{1}{2}$, der den Koeffizienten $\zeta^{m},\zeta^{\mu-m}$ vorangestellt ist, nur für den Fall, dass $\mu$ gerade ist; im anderen Fall, in dem $\mu$ ungerade ist, muss er dagegen weggelassen werden. Im letzteren Fall, in dem $\mu-m=m+1$ gilt, haben wir somit für jeden Koeffizienten $\zeta',\zeta''$ usw. einen durch $\gamma,\gamma',\gamma''$ usw. und $\delta',\delta''$ usw. bestimmten Wert; im ersteren dagegen, in dem $\mu-m=m$ gilt, für den Koeffizienten $\zeta^{\frac{1}{2}\mu}$ zwei gleiche Werte, für die auch der Wert $\delta^{\frac{1}{2}\mu}+\gamma^{\frac{1}{2}\mu}\operatorname{cotg}\mu a+\gamma^{\frac{1}{2}\mu}\operatorname{cosec}\mu a$ angenommen werden kann. Hieraus folgt
\[
\begin{aligned}
&X''={}X'-\gamma-\gamma'\cos x-\gamma''\cos 2x-\text{usw.}-\gamma^{m}\cos mx\\[5pt]
&+\operatorname{cotg}\mu a(\gamma'\sin x+\gamma''\sin 2x+\gamma'''\sin 3x+\text{usw.}+\gamma^{m}\sin mx)\\[5pt]
&+\operatorname{cosec}\mu a(\gamma\sin\mu x+\gamma'\sin(\mu-1)x\\[5pt]
&\qquad{}+\gamma''\sin(\mu-2)x+\text{usw.}+\gamma^{\mu-m}\sin(\mu-m)x).
\end{aligned}
\]
Daher ergibt sich aus $X'$ die Funktion $X''$, indem man für $\gamma$ den Ausdruck $\gamma\operatorname{cosec}\mu a\sin\mu x=\frac{\gamma\sin\mu x}{\sin\mu a}$ und für jeden Term $\gamma^{\lambda}\cos\lambda x$
\SourceLeaf{326}
\[
\begin{aligned}
&\gamma^{\lambda}\operatorname{cotg}\mu a\sin\lambda x\\[5pt]
&\qquad{}+\gamma^{\lambda}\operatorname{cosec}\mu a\sin(\mu-\lambda)x
\\[5pt]
&\text{oder}\quad
\frac{\gamma^{\lambda}\cos\mu a\sin\lambda x+\gamma^{\lambda}\sin(\mu-\lambda)x}{\sin\mu a}
\end{aligned}
\]
schreibt. Die Leser mögen diese Regel selbst mit dem vergleichen, was wir in Artikel 24 angegeben haben.

\endgroup
